\section{Exact evolution of the exterior interaction}
\label{sec:gt-interaction}

Fix \(T\ge64\), \(L=\log T\), and \(\ell=L^{-1}\). Take \(d\ge0\)
and \(\max(2,4\sqrt d)\le a\le L\), and set \(r_0=a\ell\).
On an independently enclosed count-two tube \(I=[t_0,t_1]\)
in \(0\le t\le 1/2\), consider a continued conjugation-invariant
pair of the complete \(H_t\). Its center is \(c\in[T,2T]\),
its discriminant satisfies \(|D|\le d\ell^2\), and every exterior
occurrence \(\rho_j\) has \(|\rho_j-c|\ge r_0\). The two
selected occurrences are within \(r_0/4\) of \(c\). Define
\[
 w_j=\rho_j-c,\qquad m_j=(w_j^2-D)^{-1},\qquad
 S_p=\sum_j m_j^p,\qquad M=S_1=-K .
\]
All sums count multiplicities. They are real by conjugation and
absolute by the strip and the unit count \eqref{gt-height-unit-count}.
In particular, \eqref{eq:gt-gap-law} reads \(D'=2-4MD\).

\subsection{A contour law through multiple roots}

\begin{theorem}[Interaction evolution]
\label{gt-interaction-evolution}
On the specified tube the symmetric interaction satisfies
\begin{equation}
 M'=2M^2-12S_2-16DS_3 .
 \label{gt-interaction-law}
\end{equation}
The formula holds for both signs of \(D\), a target double, and
multiple exterior roots. It also holds for a polynomial solving the
same physical heat equation on a tube with a separate counted pair.
\end{theorem}

\begin{proof}
Fix a time \(t_*\) and the circle \(\Gamma\) of radius \(r_0/2\)
about \(c(t_*)\). Its count persists near \(t_*\) by Rouché's theorem.
The source envelope \eqref{eq:tail} gives joint holomorphy in \(t,z\)
on compact sets, since the double-exponential decay dominates heat,
cosine, and derivative factors. For a field \(F\) with the same
count near \(t_*\), the contour moments
\[
 J_k=\left\langle z^k F_z/F\right\rangle_\Gamma,\qquad
 c=J_1/2,\quad D=J_2/2-c^2,\quad Q=(z-c)^2-D
\]
are locally analytic in \(t\). The quotient \(A=F/Q\) extends
jointly holomorphically and is zero-free inside the circle.
Indeed the singularities on each slice are removable by the
counted multiplicities, and the Cauchy integral of its boundary
values gives the joint extension.

Put \(g=A_z/A\), \(u=z-c\), and \(\zeta^2=D\). For the physical
heat equation \(\partial_t F=-F_{zz}\) the product rule gives
\[
 \frac{A_t}{A}=-(g_z+g^2)
      +\frac{2uc'+D'-2-4u g(c+u)}{u^2-D}.
\]
The numerator of the fraction vanishes to the multiplicities of
the roots of \(Q\). Thus the linear polynomial
\((2uc'+D'-2)/4\) is the interpolant of \(h(u)=u g(c+u)\)
at \(-\zeta,\zeta\). At \(\zeta=0\) the data are the value and
first derivative. Writing \(z_\pm=c\pm\zeta\) gives
\begin{align}
 c'&=g(z_-)+g(z_+),&
 D'&=2+4D\,g[z_-,z_+],\notag\\
 \frac{A_t}{A}&=-(g_z+g^2)-4h[-\zeta,\zeta,u].
 \label{gt-deflated-pde}
\end{align}
Here brackets denote ordinary or confluent divided differences.
The identities use analytic interpolation along the real or vertical
segment and give \(D'=2\) at a double.

By \eqref{eq:gt-regular-drift} and Proposition~\ref{gt-exterior-proposition},
the interaction is
\[
 M=-\left\langle \frac{F_z/F}{Q}\right\rangle_\Gamma .
\]
One can see the contribution \(m_j\) of each exterior pole by taking
the residue at \(\rho_j\) of \(1/((z-\rho_j)Q(z))\), whose residue
at infinity is zero. The term \(Q_z/Q^2=-(1/Q)_z\) from the two
target poles integrates to zero. The formula also applies at a
target double.

First take a polynomial slice on which all roots are simple.
At a root \(\rho\), implicit differentiation gives
\(\rho'=F_{zz}(\rho)/F_z(\rho)=2\sum_{\eta\ne \rho}
(\rho-\eta)^{-1}\), with \(\eta\) ranging over the other roots.
Write \(P=\sum_j w_jm_j\), \(V=\sum_j w_jm_j^2\), with only
exterior indices. The first identity of \eqref{gt-deflated-pde}
gives \(c'=-2P\), and the two target contributions to each exterior
velocity are \(4w_jm_j\). Consequently
\[
 w_j'=4w_jm_j+2P+2\sum_{k\ne j}(w_j-w_k)^{-1}.
\]
Differentiating \(m_j\), using \(w_j^2m_j=1+Dm_j\), and
pairing the \(j,k\) and \(k,j\) terms yields
\begin{align*}
 M'={}&(2-4DM)S_2-8(S_2+DS_3)-4PV\\
       &\quad-4\sum_{j<k}
          \frac{w_jm_j^2-w_km_k^2}{w_j-w_k}.
\end{align*}
Subtracting \(w_jm_j-w_km_k\) first gives
\((w_jm_j-w_km_k)/(w_j-w_k)=-(D+w_jw_k)m_jm_k\).
Multiplying this quotient by \(m_j+m_k\) and subtracting \(m_jm_k\)
gives
\[
 \frac{w_jm_j^2-w_km_k^2}{w_j-w_k}
 =-m_jm_k-(D+w_jw_k)(m_j^2m_k+m_jm_k^2).
\]
In the sum over \(j<k\), the three products with factors \(1,D,w_jw_k\)
on the right, before the minus signs, give respectively
\[
 \frac{M^2-S_2}{2},\qquad D(MS_2-S_3),\qquad PV-S_2-DS_3 .
\]
The last diagonal is \(\sum_j w_j^2m_j^3=S_2+DS_3\).
Substitution cancels \(PV\) and \(DMS_2\) and gives
\eqref{gt-interaction-law}.

At a multiple-root polynomial slice, perturb its finitely many
root factors to distinct locations while preserving the count on
each side of \(\Gamma\). Give each perturbed slice its polynomial
heat evolution. The contour formulas at the fixed time give
\begin{align}
 \partial_t(F_z/F)&=-F_{zzz}/F+F_zF_{zz}/F^2,\notag\\
 M'&=-\left\langle
      \frac{\partial_t(F_z/F)}{Q}
        -\frac{(F_z/F)Q_t}{Q^2}
      \right\rangle_\Gamma,\qquad
 Q_t=-2(z-c)c'-D'.
 \label{gt-interaction-contour-derivative}
\end{align}
The same first expression gives \(c',D'\) by differentiating the
moments. The lower bounds for the denominators on \(\Gamma\)
persist under the perturbation. Thus \(M'\) passes to the original
slice, as do \(M,S_2,S_3\), since the exterior denominators there
stay separated from zero. This proves the finite identity through
both target and exterior multiplicities.

For \(H_t\), the order argument preceding
Proposition~\ref{gt-exterior-proposition} applies uniformly through
\(t=1/2\) using \(u^2/2\) in its integrand bound. Also \(H_t(0)>0\).
The strip and unit count give \(\sum_\rho|\rho|^{-2}<\infty\) by
unit intervals on both tails. In the genus-one product at \(t_*\),
opposite factors combine to \(1-z^2/\rho^2\). Their locally
uniformly convergent product is even, so evenness of \(H_{t_*}\)
forces the linear exponential in the Hadamard prefactor to be
constant. Truncate this product at increasing moduli, with all
opposite and conjugate factors and their multiplicities, keeping the
constant prefactor \(H_{t_*}(0)\). Each truncation is the time slice
of a polynomial heat solution, obtained by applying the terminating
operator \(\exp(-(t-t_*)\partial_z^2)\).

On the single slice \(t_*\), compact product convergence and
Cauchy's formula give convergence of every spatial derivative,
including on \(\Gamma\). The time derivatives of the logarithmic
derivatives then converge by
\eqref{gt-interaction-contour-derivative}. The contour moments give
convergence of \(c,D,c',D'\), and the same contour derivative gives
convergence of \(M'\). Once the target factors are present, each
retained exterior summand at \(t_*\) already uses the canonical
\(c,D\). Absolute convergence gives the limits of \(S_2,S_3\).
Taking the limit proves \eqref{gt-interaction-law} at each \(t_*\).
The polynomial time neighborhoods may depend on the truncation.
\end{proof}

\subsection{Complete shell estimates and Riccati comparison}

Use the fixed normalization
\(\sigma=(t-t_0)/\ell^2\), \(m=\ell^2 M\),
\(\mu_j=\ell^2m_j\), and \(\delta=D/\ell^2\).
Set \(q_0=16/15\). The collar gives
\(|D|\le |w_j|^2/16\), hence
\(|\mu_j|\le q_0 L^{-2}|w_j|^{-2}\). Let \(C_{\rm u}\ge1\)
bound both the constant in \eqref{gt-height-unit-count} and
the safe-line constant in \eqref{gt-safe-line}.
Within \(|\Re w_j|<1\) the two closed unit intervals give at most
\(4C_{\rm u}L\) occurrences, each with
\(|\mu_j|\le q_0a^{-2}\). For \(n\ge1\) the shell
\(n\le|\Re w_j|<n+1\) lies in unit intervals starting at
\(c-n-1,c+n\). Both starts \(v\) satisfy
\(2+|v|\le 3T(n+1)\), giving shell count at most
\(2C_{\rm u}(2L+\log(n+1))\). Each summand there is bounded by
\(|\mu_j|\le q_0 L^{-2}n^{-2}\). The integral bound for
\(\sum_{n\ge1}n^{-2}\) is \(2\). Also \(\log(n+1)\le\sqrt n\),
so \(\sum_{n\ge1}n^{-2p}\log(n+1)\le3\) for \(p\ge1\), by the
integral bound for \(\sum_{n\ge1}n^{-3/2}\). Thus the shells with
\(n\ge1\) contribute at most
\(2C_{\rm u}q_0^p L^{-2p}(4L+3)\le10C_{\rm u}q_0^p L^{1-2p}\).
Combining with the near band and using \(a\le L\) gives
\begin{equation}
 \sum_j|\mu_j|^p\le B_pL,\qquad
 B_p=\frac{14C_{\rm u}q_0^p}{a^{2p}}\quad(p=1,2,3).
 \label{gt-interaction-shell}
\end{equation}

The same shell counts also bound the deflated logarithmic derivative
on \(|z-c|\le r_0/2\). At \(c+4\ii\), removing the target terms
from \eqref{gt-safe-line} gives \(|g(c+4\ii)|\le2C_{\rm u}L+1\).
Subtracting \(g(c+4\ii)\) cancels the regularizing constants in the
logarithmic derivative and leaves the differences
\((z-\rho_j)^{-1}-(c+4\ii-\rho_j)^{-1}\).
The near differences are bounded by \(2/r_0+1/3\), using the
strip and the collar. On shell \(n\ge1\) the two denominator
factors are at least \(n/2,n\) and the numerator at most \(9/2\),
so each difference is at most \(9/n^2\). The preceding sums give
\[
 |g(z)|\le 8C_{\rm u}L^2/a+95C_{\rm u}L
 \qquad (|z-c|\le r_0/2).
\]
The exterior reciprocal-power sums for positive spatial derivatives
also converge uniformly on that disk by the same shell counts.

In the scaled time, \eqref{gt-interaction-law} becomes
\begin{equation}
 \frac{\dd m}{\dd\sigma}=2m^2-\mathcal F,\qquad
 \mathcal F=12\sum_j\mu_j^2+16\delta\sum_j\mu_j^3 .
 \label{gt-interaction-scaled}
\end{equation}
Write \(B=B_1\) and \(C_F=12B_2+16dB_3>0\).
Then \(|m|\le BL\), \(|\mathcal F|\le C_FL\), and
\(\dd m/\dd\sigma\ge 2m^2-C_FL\).

\begin{corollary}[Lower barrier and integrated adverse interaction]
\label{gt-riccati-barrier}
Put \(h=\sqrt{C_FL/2}\), \(k=\max(BL/h,2)\), and
\(u_0=\operatorname{arccoth}k\). Within the tube,
\begin{align}
 m(\sigma)&\ge-h\coth(2h\sigma+u_0)
              \ge-h\coth(2h\sigma)\qquad(\sigma>0),
 \label{gt-riccati-lower}\\
 \int_0^s(-m)_+\,\dd\sigma
    &\le \frac12\log\frac{\sinh(2hs+u_0)}{\sinh u_0}
       \le hs+\frac12\log\frac{k+1}{2}.
 \label{gt-riccati-cost}
\end{align}
\end{corollary}
\begin{proof}
The function \(v(\sigma)=-h\coth(2h\sigma+u_0)\) solves
\(v'=2v^2-C_FL\) and \(v(0)=-hk\le m(0)\). The inequality
\((m-v)'-2(m+v)(m-v)\ge0\) and its positive integrating factor
imply \(m\ge v\). The second lower bound follows from monotonicity
of \(\coth\). Integrating \(-v\) gives the logarithm in
\eqref{gt-riccati-cost}. Bounding the numerator \(\sinh\) by half
its positive exponential and using
\(\e^{u_0}/(2\sinh u_0)=(k+1)/2\) gives the second inequality.
\end{proof}

At any fixed positive elapsed scaled-time buffer, the lower allowance
is of order \(-\sqrt L\) for fixed \(a,d,C_{\rm u}\). Over a
scaled interval of length \(s\) the multiplier in
\eqref{eq:gt-backward-gap} is bounded by
\begin{equation}
 \exp\!\left(-4\int_{t_0}^{t_0+\ell^2s} M\,\dd t\right)
 \le \e^{4hs}\left(\frac{k+1}{2}\right)^2
 =O_{a,d,C_{\rm u}}(L)\exp\!\big(O_{a,d,C_{\rm u}}(s\sqrt L)\big).
 \label{gt-riccati-multiplier}
\end{equation}
This uses \(\int M\,\dd t=\int m\,\dd\sigma\). Restarting the
comparison at any interior time gives the same estimate for each
kernel in the signed cost integral in \eqref{eq:gt-backward-gap}.
The pointwise count alone gives \(\e^{4BLs}\). The improved bound
still permits fixed negative normalized margins and growth of the
multiplier at a fixed positive \(s\). Applying the gap equation
also requires a separately enclosed value of \(D\) at the chosen
anchor. The initial count bound used in the comparison concerns \(m\).

\subsection{Negative visibility and conversion to a slice count}

Write \(C_1=2B^2+C_F\). By \eqref{gt-interaction-scaled},
\[
 \left|\frac{\dd}{\dd t}(\ell^2M)\right|
 \le L^2(2B^2L^2+C_FL)\le C_1L^4 .
\]
For a slice margin \(\ell^2M(t_*)\le-\epsilon\), \(\epsilon>0\),
this guarantees \(\ell^2M\le-\epsilon/2\) on a chosen side of length
\(\Delta=\epsilon/(2C_1L^4)\) when that interval lies in the
independently enclosed tube. A strict sign without a margin allows
the guaranteed length to shrink to zero.

For \(n_\epsilon\) distinct qualifying pairs at a snapshot, take
a common one-sided interval of length \(\Delta>0\) within all of
their paid tubes on which each interaction stays negative. The
continued pairs remain distinct there. If two coincide at a time,
the same local counting contour fixes the unique continuation near
that time. Coincidence is also closed by continuity and would then
extend to the snapshot. Thus the \(n_\epsilon\) pairs contribute at
least \(\Delta n_\epsilon\) to the negative pair-time occupation.
By Theorem~\ref{gt-pair-occupation} applied on the common interval,
\begin{equation}
 n_\epsilon\le C_{a,b}\frac{T}{L\Delta}
       +\frac{32D_\chi}{3}\left(\frac12+\frac{L}{3a}\right)L^3,
 \qquad C_{a,b}=B_b\left(\frac{93}{b^2}+\frac{279}{ab^3}\right),
 \label{gt-visibility-comparison}
\end{equation}
with the density and cutoff constants of that theorem. The
modulus-guaranteed interval gives \(O(TL^3/\epsilon)+O(L^4)\)
for fixed \(a,b,d\) and source constants. At fixed \(\epsilon\)
its leading order exceeds the slice estimate in
Theorem~\ref{gt-pair-negative-count} by two powers of \(L\).
Even an independently guaranteed negative interval of length
comparable to \(\ell^2\) gives only \(O(TL)+O(L^4)\).
For every \(\Delta\le1/2\) the first term of the explicit allowance
is at least \(2C_{a,b}T/L\). Hence this conversion cannot furnish an
upper allowance below one at growing \(T\) with fixed collar parameter
\(a\). These comparisons are between upper allowances for the
specified pair class. They give no lower bound on actual exceptions
or their lifetimes.

The finite quadratic and even quartic heat fields obey
\eqref{gt-interaction-law}, including on collars through a collision.
Their finite divisors cannot supply the canonical growing population,
though they can satisfy coarse unit upper counts. The canonical
comparison above uses the complete source through the unit count
and occupation estimates. A direct calculation with the same
signed interaction, tube, source inputs, and independent anchor
has the same comparison and the same visibility requirements
\cite{gtdynamics}.
